\section{Conclusion}
\label{sec:conclusion}

GRASP provides an evidence-linked workflow for assessing disclosure characteristics and greenwashing risk in structured Indian BRSR reports. The central contribution is the connection between report-derived principle labels, language and numerical review checks, and company scores with visible component calculations. Shared report identifiers allow an analyst to inspect the text or filed value behind a result, while company and sector views support review at different levels. The vagueness detector also supplies the Clarity component of the Transparency Score, maintaining a direct link between text screening and company-level scoring.

The recovered classifier achieves 93.00\% Top-1 agreement and macro-F1 of 0.9158 against weak labels on 6,201 text segments from 175 companies excluded from supervised training. Source-wise and excluded-concept results provide further context for the overall agreement. The fixed language detector achieves F1 of 0.7436 on 80 development text segments, and the Transparency Score correlates with 20 historical human ratings at $r=0.6314$. Application outputs cover 1,164 companies in 72 recorded sectors. The four reporting levels show how component differences explain company and sector scores, while the eight paired-year cases record a mean increase of 5.23 points. Each result remains tied to the evaluation reference and company group described in the paper.

This work can be extended using multiple levels. At the data level, collection can include annexures, charts, and additional annual reports. At the modelling level, larger labelled datasets, blind human-rating studies, and adaptation restricted to declared training companies can support further evaluation on newly collected companies. Operational business descriptions can refine sector assignment and reporting-item applicability. Numerical checks can also compare measure definitions, units, periods, and business boundaries before generating review records. As more reporting years become available, repeated company comparisons can form a longer panel under a recorded score version. Finally, the evidence-linking approach can be evaluated against international sustainability frameworks using separate reporting-item mappings and tests for each standard.
